\documentclass[conference]{IEEEtran}
\IEEEoverridecommandlockouts
\usepackage{cite}
\usepackage{amsmath,amssymb,amsfonts}
\usepackage{algorithmic}
\usepackage{graphicx}
\usepackage{textcomp}
\usepackage{xcolor}
\usepackage[utf8]{inputenc}
\usepackage[T1]{fontenc}
\def\BibTeX{{\rm B\kern-.05em{\sc i\kern-.025em b}\kern-.08em
    T\kern-.1667em\lower.7ex\hbox{E}\kern-.125emX}}
\begin{document}

\title{Avaliação de Características Hierárquicas de Redes Pré-treinadas para Recuperação de Padrões de Cor e Textura
\\
\thanks{}
}

\author{
\IEEEauthorblockN{Leonardo [Sobrenome]}
\IEEEauthorblockA{\textit{Departamento de Ciência da Computação} \\
\textit{Universidade Tecnológica Federal do Paraná (UTFPR)}\\
Santa Helena (PR), Brasil \\
email@alunos.utfpr.edu.br}
\and
\IEEEauthorblockN{Alceu de Souza Britto Jr.}
\IEEEauthorblockA{\textit{Departamento de Ciência da Computação} \\
\textit{Pontifícia Universidade Católica do Paraná (PUCPR)}\\
Curitiba (PR), Brasil \\
alceu@ppgia.pucpr.br}

\and\and

\hspace{1cm}
\IEEEauthorblockN{[Nome do Amigo]}
\hspace{9cm}\IEEEauthorblockA{\textit{Departamento de Ciência da Computação} \\
\textit{Universidade Tecnológica Federal do Paraná (UTFPR)}\\
\hspace{1cm}Santa Helena (PR), Brasil \\
\hspace{1cm}email@alunos.utfpr.edu.br}
\and

\IEEEauthorblockN{Arlete Teresinha Beuren}
\IEEEauthorblockA{\textit{Departamento de Ciência da Computação} \\
\textit{Universidade Tecnológica Federal do Paraná (UTFPR)}\\
Santa Helena (PR), Brasil \\
arletebeuren@utfpr.edu.br}
}

\maketitle

\begin{abstract}
Este estudo investiga as capacidades de extração hierárquica de características de arquiteturas modernas de redes neurais pré-treinadas para reconhecimento de atributos visuais, especificamente direcionado a atributos de cor e textura em imagens de roupas. Avaliamos quatro arquiteturas backbone distintas representando diferentes paradigmas: VGG-16 e ResNet-50 como redes neurais convolucionais clássicas, iBOT como um Vision Transformer auto-supervisionado, e VMamba como um Modelo Visual de Espaço de Estados. Características são extraídas de camadas intermediárias de cada arquitetura e avaliadas usando classificação k-Nearest Neighbors com validação cruzada de 5 folds. Nossos experimentos revelam que diferentes profundidades arquiteturais exibem sensibilidades variadas a atributos de cor e textura, com camadas iniciais consistentemente capturando informação de cor enquanto camadas mais profundas melhor representam padrões de textura. Os resultados demonstram que o iBOT alcança o melhor reconhecimento de textura (97,2\% no bloco 9) enquanto ResNet-50 e iBOT empatam para reconhecimento de cor (95,9\%). Estes achados fornecem insights sobre a especialização de características por camada através de diferentes paradigmas de redes neurais e orientação prática para sistemas de recuperação de imagens baseados em atributos.

\end{abstract}

\begin{IEEEkeywords}
Redes Neurais Convolucionais, Vision Transformers, Modelos Visuais de Espaço de Estados, Extração de Características, Atributos
\end{IEEEkeywords}

\section{Introdução}
Os avanços em aprendizado profundo têm possibilitado soluções de visão computacional cada vez mais sofisticadas para aplicações industriais. Um exemplo notável é a recuperação de imagens em e-commerce, onde buscas podem ser conduzidas usando texto, imagens ou uma combinação de ambos—componentes-chave de sistemas inovadores. Enquanto a busca baseada em texto permanece como o método de consulta mais comum, sua eficácia depende fortemente de descrições detalhadas dos produtos. Para melhorar a experiência do usuário, alguns motores de busca incorporam busca baseada em imagem, oferecendo resultados significativamente mais precisos [1].

Motores de busca baseados em conteúdo de imagem apresentam desafios únicos e encontram aplicações além do e-commerce, incluindo recuperação de imagens médicas, sistemas de informação geográfica e vigilância [12]. O principal desafio está em extrair com precisão o conteúdo relevante de uma imagem, particularmente quando o interesse do usuário é baseado em atributos específicos como cor, textura ou sua combinação. No contexto do e-commerce, este tipo de busca focada em atributos está cada vez mais em demanda.

Pesquisas anteriores demonstraram que diferentes camadas de redes neurais convolucionais (CNNs) capturam diferentes tipos de informação visual [2]. Camadas iniciais tipicamente capturam características de baixo nível como bordas e cores, enquanto camadas mais profundas codificam informação mais abstrata e semântica, incluindo texturas e padrões. Esta representação hierárquica de características sugere que selecionar camadas apropriadas para atributos visuais específicos poderia melhorar significativamente a acurácia de reconhecimento.

Avanços recentes em arquiteturas de redes neurais introduziram novos paradigmas além das CNNs tradicionais. Vision Transformers (ViTs) aplicam a arquitetura transformer ao processamento de imagens tratando imagens como sequências de patches [3]. Mais recentemente, Modelos Visuais de Espaço de Estados como VMamba emergiram, combinando a eficiência de modelos de espaço de estados com extração de características visuais [4]. Entender como estes diferentes paradigmas arquiteturais organizam características hierárquicas é crucial para reconhecimento efetivo de atributos visuais.

Este estudo avalia quatro arquiteturas backbone distintas para extração de características de cor e textura de imagens de roupas: VGG-16 [5], uma CNN clássica com camadas convolucionais e de pooling sequenciais; ResNet-50 [6], uma CNN profunda com conexões residuais; iBOT [7], um Vision Transformer auto-supervisionado pré-treinado usando modelagem de imagem mascarada; e VMamba [4], um Modelo Visual de Espaço de Estados que aplica a arquitetura Mamba a tarefas visuais.

Para guiar nossos experimentos, formulamos as seguintes questões de pesquisa:
(QP1) Quais camadas dentro de cada arquitetura fornecem a melhor representação de atributos de cor?
(QP2) Quais camadas dentro de cada arquitetura fornecem a melhor representação de atributos de textura?
(QP3) Como os quatro paradigmas arquiteturais se comparam em sua organização hierárquica de características?

Este estudo está organizado da seguinte forma: a Seção 2 revisa trabalhos relacionados, a Seção 3 descreve as arquiteturas backbone, a Seção 4 detalha a metodologia aplicada, a Seção 5 apresenta resultados experimentais, e a Seção 6 conclui com achados e direções futuras.


\section{TRABALHOS RELACIONADOS}

O uso de redes neurais profundas para similaridade visual e reconhecimento de atributos tem sido extensivamente estudado na literatura. Motores de busca incorporam cada vez mais consultas baseadas em imagem, avançando além das buscas tradicionais baseadas em texto [1].

Baloian et al. [2] investigaram camadas da ResNet-50 para determinar representações ótimas para buscas de cor e textura. Sua pesquisa concluiu que camadas iniciais da rede são mais adequadas para reconhecimento de cor, enquanto camadas mais profundas se destacam no reconhecimento de textura. Características foram extraídas de blocos residuais usando Global Average Pooling (GAP), e um classificador k-Nearest Neighbors (k-NN) com validação cruzada de 5 folds foi empregado para avaliação.

Os autores em [8] classificam espécies de plantas por imagem de folha usando a arquitetura de redes neurais Siamesas, combinando características de camadas convolucionais intermediárias para melhorar representações. Segundo os autores, a combinação de características de diferentes camadas apresenta um ganho de desempenho relevante e diferentes camadas trazem resultados diferentes.

No estudo de [9], os autores usaram uma rede neural Siamesa para definir o índice de similaridade das imagens de cada classe através de uma análise da folha inteira e partes dela.

Em [10], os autores apresentam um método eficiente de extração de características baseado em pacotes wavelet e autovalores de filtros de Gabor para análise de textura, combinado com k-NN para classificação. Estudos também propõem a combinação de características de cor e textura como vetor de características [11]. Lin et al. [13] propuseram características de conteúdo de imagem baseadas em cor, textura e distribuição de cor usando matrizes de co-ocorrência de cor e métodos baseados em histograma.

Avanços recentes em aprendizado auto-supervisionado introduziram estratégias poderosas de pré-treinamento para representações visuais. Vision Transformers pré-treinados com métodos como DINO [14] e iBOT [7] demonstraram capacidades notáveis de extração de características sem requerer dados rotulados durante o pré-treinamento. O surgimento de Modelos de Espaço de Estados para visão, particularmente VMamba [4], oferece um novo paradigma combinando eficiência com extração hierárquica de características visuais.


\section{ARQUITETURAS BACKBONE}

Esta seção descreve as quatro arquiteturas backbone avaliadas neste estudo.

\subsection{VGG-16}
A VGG-16 [5] é uma rede neural convolucional clássica caracterizada por sua arquitetura uniforme usando pequenos filtros convolucionais (3×3) em toda a rede. A arquitetura consiste em 16 camadas com pesos organizadas em 5 blocos convolucionais, cada um seguido por max-pooling. Características são extraídas após cada camada de max-pooling usando Global Average Pooling (GAP), resultando em dimensões de 64, 128, 256, 512 e 512 para as camadas 1 a 5.

\subsection{ResNet-50}
A ResNet-50 [6] introduziu conexões residuais (skip connections) que permitem o treinamento de redes muito mais profundas. A arquitetura consiste em 4 blocos residuais principais. Características são extraídas de cada bloco usando GAP, com dimensões de 256, 512, 1024 e 2048 para as camadas 1 a 4.

\subsection{iBOT}
O iBOT (Image BERT Pre-Training with Online Tokenizer) [7] é um Vision Transformer pré-treinado usando modelagem de imagem mascarada auto-supervisionada. A variante ViT-S/16 consiste em 12 blocos transformer, cada um produzindo embeddings de 384 dimensões. Características são extraídas usando o token CLS de cada bloco.

\subsection{VMamba}
O VMamba [4] adapta a arquitetura Mamba—um modelo seletivo de espaço de estados—para tarefas visuais. Diferente de transformers com complexidade quadrática, o VMamba alcança complexidade linear. A variante VMamba-Tiny tem 4 estágios com dimensões de 96, 192, 384 e 768.


\section{METODOLOGIA APLICADA}

A metodologia consiste em extrair características de camadas intermediárias de redes backbone pré-treinadas e avaliar sua eficácia para reconhecimento de cor e textura usando um classificador k-NN.

\subsection{Conjunto de Dados}
O conjunto de dados usado nesta pesquisa foi criado por Baloian et al. [2], utilizando imagens obtidas do Kaggle e outros repositórios online. Apresenta uma coleção diversa de itens de vestuário incluindo chapéus, vestidos, camisas, meias, bolsas e cachecóis.

Conjunto de Dados de Cor: Contém 1.000 imagens organizadas em 10 classes (100 imagens cada): vermelho, preto, azul, verde, amarelo, cinza, marrom, rosa, roxo e laranja.

Conjunto de Dados de Textura: Contém 1.000 imagens organizadas em 10 classes (100 imagens cada): quadrados, listras, flores, leopardo, bolinhas, básico, paisley, argyle, pata de galo e lantejoulas.

\begin{figure}[htbp]
\centerline{\includegraphics[width=0.48\textwidth]{fig_texture_samples.png}}
\caption{Imagens de amostra do conjunto de dados de textura mostrando seis classes de textura: Argyle, Básico, Pata de Galo, Flores, Leopardo e Paisley.}
\label{fig_texture}
\end{figure}

\begin{figure}[htbp]
\centerline{\includegraphics[width=0.48\textwidth]{fig_color_samples.png}}
\caption{Imagens de amostra do conjunto de dados de cor mostrando todas as dez classes de cor: Amarelo, Azul, Marrom, Cinza, Roxo, Laranja, Preto, Vermelho, Rosa e Verde.}
\label{fig_color}
\end{figure}

O banco de dados foi dividido em 70\% para treinamento e 30\% para teste.

\subsection{Processo de Extração de Características}
Para cada arquitetura backbone, características são extraídas de camadas intermediárias da seguinte forma: (1) Imagens são redimensionadas para 224×224 pixels e normalizadas usando estatísticas do ImageNet; (2) Hooks de forward capturam ativações intermediárias durante a inferência; (3) Global Average Pooling (GAP) é aplicado para CNNs e VMamba, enquanto extração do token CLS é usada para iBOT; (4) Características extraídas são armazenadas para classificação.

\subsection{Protocolo de Classificação}
O algoritmo k-Nearest Neighbors é usado com k=5 vizinhos, métrica de distância Euclidiana e validação cruzada estratificada de 5 folds.


\section{RESULTADOS EXPERIMENTAIS}

Esta seção apresenta os resultados experimentais para cada arquitetura backbone.

\subsection{Resultados VGG-16}

\begin{table}[htbp]
\caption{Acurácia por Camada da VGG-16 para Cor e Textura.}
\begin{center}
\begin{tabular}{|c|c|c|c|}
\hline
\textbf{Camada} & \textbf{Cor(\%)} & \textbf{Textura(\%)} & \textbf{Dim}\\
\hline
VGG-16\_Camada\_1 & \textbf{95,1} & 70,7 & 64 \\
\hline
VGG-16\_Camada\_2 & 90,3 & 86,6 & 128 \\
\hline
VGG-16\_Camada\_3 & 85,5 & 93,2 & 256 \\
\hline
VGG-16\_Camada\_4 & 71,8 & \textbf{94,7} & 512 \\
\hline
VGG-16\_Camada\_5 & 55,1 & 89,4 & 512 \\
\hline
\end{tabular}
\label{tab1}
\end{center}
\end{table}

A VGG-16 demonstra clara especialização por camada: a Camada 1 alcança a maior acurácia de cor (95,1\%), enquanto a Camada 4 alcança a melhor acurácia de textura (94,7\%).

\subsection{Resultados ResNet-50}

\begin{table}[htbp]
\caption{Acurácia por Camada da ResNet-50 para Cor e Textura.}
\begin{center}
\begin{tabular}{|c|c|c|c|}
\hline
\textbf{Camada} & \textbf{Cor(\%)} & \textbf{Textura(\%)} & \textbf{Dim}\\
\hline
ResNet-50\_Camada\_1 & 95,5 & 83,4 & 256 \\
\hline
ResNet-50\_Camada\_2 & \textbf{95,9} & 90,7 & 512 \\
\hline
ResNet-50\_Camada\_3 & 88,3 & \textbf{93,3} & 1024 \\
\hline
ResNet-50\_Camada\_4 & 59,2 & 86,3 & 2048 \\
\hline
\end{tabular}
\label{tab2}
\end{center}
\end{table}

A ResNet-50 mostra padrões similares: a Camada 2 alcança a melhor acurácia de cor (95,9\%), enquanto a Camada 3 alcança a melhor acurácia de textura (93,3\%).

\subsection{Resultados iBOT}

\begin{table}[htbp]
\caption{Acurácia por Bloco do iBOT para Cor e Textura.}
\begin{center}
\begin{tabular}{|c|c|c|c|}
\hline
\textbf{Bloco} & \textbf{Cor(\%)} & \textbf{Textura(\%)} & \textbf{Dim}\\
\hline
iBOT\_Bloco\_0 & 94,5 & 52,2 & 384 \\
\hline
iBOT\_Bloco\_2 & \textbf{95,9} & 79,9 & 384 \\
\hline
iBOT\_Bloco\_5 & 88,7 & 91,0 & 384 \\
\hline
iBOT\_Bloco\_9 & 69,9 & \textbf{97,2} & 384 \\
\hline
iBOT\_Bloco\_11 & 62,8 & 96,0 & 384 \\
\hline
\end{tabular}
\label{tab3}
\end{center}
\end{table}

O iBOT exibe um trade-off pronunciado entre cor e textura. O Bloco 2 alcança a melhor cor (95,9\%), enquanto o Bloco 9 alcança a maior acurácia de textura (97,2\%).

\subsection{Resultados VMamba}

\begin{table}[htbp]
\caption{Acurácia por Estágio do VMamba para Cor e Textura.}
\begin{center}
\begin{tabular}{|c|c|c|c|}
\hline
\textbf{Estágio} & \textbf{Cor(\%)} & \textbf{Textura(\%)} & \textbf{Dim}\\
\hline
VMamba\_Estágio\_1 & \textbf{94,4} & 93,9 & 96 \\
\hline
VMamba\_Estágio\_2 & 93,4 & \textbf{95,8} & 192 \\
\hline
VMamba\_Estágio\_3 & 78,3 & \textbf{95,8} & 384 \\
\hline
VMamba\_Estágio\_4 & 54,1 & 86,5 & 768 \\
\hline
\end{tabular}
\label{tab4}
\end{center}
\end{table}

O VMamba mostra Estágio 1 ótimo para cor (94,4\%) e Estágios 2-3 para textura (95,8\%).

\subsection{Análise Comparativa}

\begin{table}[htbp]
\caption{Resumo dos Melhores Resultados Entre Arquiteturas.}
\begin{center}
\begin{tabular}{|c|c|c|c|}
\hline
\textbf{Modelo} & \textbf{Atributo} & \textbf{Melhor Camada} & \textbf{Acc(\%)}\\
\hline
VGG-16 & Cor & Camada 1 & 95,1 \\
\hline
VGG-16 & Textura & Camada 4 & 94,7 \\
\hline
ResNet-50 & Cor & Camada 2 & 95,9 \\
\hline
ResNet-50 & Textura & Camada 3 & 93,3 \\
\hline
iBOT & Cor & Bloco 2 & 95,9 \\
\hline
iBOT & Textura & Bloco 9 & \textbf{97,2} \\
\hline
VMamba & Cor & Estágio 1 & 94,4 \\
\hline
VMamba & Textura & Estágio 2-3 & 95,8 \\
\hline
\end{tabular}
\label{tab5}
\end{center}
\end{table}

Principais achados: (1) ResNet-50 e iBOT empatam para cor em 95,9\%; (2) iBOT alcança a maior acurácia de textura em 97,2\%; (3) Todas as arquiteturas mostram camadas iniciais ótimas para cor e camadas mais profundas para textura.

\subsection{DISCUSSÃO}

Resultados promissores foram observados em todas as quatro arquiteturas. O achado mais importante é que, independentemente do paradigma arquitetural (CNN, Transformer ou Modelo de Espaço de Estados), a tendência é mantida: camadas iniciais apresentam melhores resultados para cor e camadas mais profundas para textura.

\begin{figure}[htbp]
\centerline{\includegraphics[width=0.48\textwidth]{fig_texture_results.png}}
\caption{Exemplos de classificação de textura. Rótulos superiores mostram ground truth, rótulos inferiores mostram predições. O modelo identifica corretamente padrões desafiadores como leopardo e paisley.}
\label{fig_texture_results}
\end{figure}

\begin{figure}[htbp]
\centerline{\includegraphics[width=0.48\textwidth]{fig_color_results.png}}
\caption{Exemplos de classificação de cor. Rótulos superiores mostram ground truth, rótulos inferiores mostram predições. Alguma confusão ocorre entre cores similares (vermelho/laranja, laranja/marrom).}
\label{fig_color_results}
\end{figure}

O pré-treinamento auto-supervisionado do iBOT produz características hierárquicas fortes, alcançando o melhor reconhecimento de textura (97,2\%) enquanto mantém reconhecimento de cor competitivo (95,9\%). A VGG-16 alcança alto desempenho com características de baixa dimensionalidade (64D para cor), sugerindo que representações eficientes são possíveis com arquiteturas simples.


\section{CONCLUSÕES}

Este estudo avaliou a extração hierárquica de características de quatro arquiteturas de redes neurais pré-treinadas—VGG-16, ResNet-50, iBOT e VMamba—para reconhecimento de atributos visuais em imagens de roupas.

Os experimentos demonstram que: (1) Todas as quatro arquiteturas exibem especialização hierárquica de características, com camadas iniciais ótimas para cor e camadas mais profundas para textura; (2) O iBOT alcança o melhor reconhecimento de textura (97,2\%), enquanto ResNet-50 e iBOT empatam para cor (95,9\%); (3) A organização hierárquica é consistente através dos paradigmas CNN, Transformer e SSM.

Para experimentos futuros, recomendamos investigar estratégias de fusão combinando características de múltiplas camadas para aproveitar informações de cor e textura simultaneamente, bem como explorar abordagens de fusão aprendida que poderiam eliminar a necessidade de seleção manual de camadas.


\section*{Agradecimentos}
Este trabalho é apoiado pelo Programa de Pós-Graduação em Informática da Pontifícia Universidade Católica do Paraná e parcialmente apoiado pelo Conselho Nacional de Desenvolvimento Científico e Tecnológico (CNPq) vinculado ao Ministério da Ciência, Tecnologia e Inovação.


\begin{thebibliography}{00}
\bibitem{b1} Dubey, S. R. (2021). A decade survey of content based image retrieval using deep learning. IEEE Transactions on Circuits and Systems for Video Technology, 32(5), 2687-2704.
\bibitem{b2} Baloian, A., Murrugarra-Llerena, N., \& Saavedra, J. M. (2021). Scalable visual attribute extraction through hidden layers of a residual convnet. arXiv preprint arXiv:2104.00161.
\bibitem{b3} Dosovitskiy, A. et al. (2020). An image is worth 16x16 words: Transformers for image recognition at scale. arXiv preprint arXiv:2010.11929.
\bibitem{b4} Liu, Y. et al. (2024). VMamba: Visual state space model. arXiv preprint arXiv:2401.10166.
\bibitem{b5} Simonyan, K., \& Zisserman, A. (2014). Very deep convolutional networks for large-scale image recognition. arXiv preprint arXiv:1409.1556.
\bibitem{b6} He, K., Zhang, X., Ren, S., \& Sun, J. (2016). Deep residual learning for image recognition. In CVPR (pp. 770-778).
\bibitem{b7} Zhou, J. et al. (2021). iBOT: Image BERT pre-training with online tokenizer. arXiv preprint arXiv:2111.07832.
\bibitem{b8} Moresco, M., Britto, A. D. S., Costa, Y. M., Senger, L. J., \& Hochuli, A. G. (2022). Combining multi-layer features for plant species classification in a Siamese network. In IEEE SMC (pp. 2446-2451).
\bibitem{b9} Araújo, V. M., Britto Jr, A. S., Oliveira, L. S., \& Koerich, A. L. (2022). Two-view fine-grained classification of plant species. Neurocomputing, 467, 427-441.
\bibitem{b10} Irtaza, A., Jaffar, M. A., Aleisa, E., \& Choi, T. S. (2014). Embedding neural networks for semantic association in content based image retrieval. Multimedia Tools and Applications, 72, 1911-1931.
\bibitem{b11} Atlam, H. F., Attiya, G., \& El-Fishawy, N. (2017). Integration of color and texture features in CBIR system. Int. J. Comput. Appl., 164(3), 23-29.
\bibitem{b12} Youssef, S. M. (2012). ICTEDCT-CBIR: Integrating curvelet transform with enhanced dominant colors extraction and texture analysis for efficient content-based image retrieval. Computers \& Electrical Engineering, 38(5), 1358-1376.
\bibitem{b13} Lin, C. H., Chen, R. T., \& Chan, Y. K. (2009). A smart content-based image retrieval system based on color and texture feature. Image and Vision Computing, 27(6), 658-665.
\bibitem{b14} Caron, M. et al. (2021). Emerging properties in self-supervised vision transformers. In ICCV (pp. 9650-9660).
\end{thebibliography}

\end{document}
